% Options for packages loaded elsewhere
\PassOptionsToPackage{unicode}{hyperref}
\PassOptionsToPackage{hyphens}{url}
\PassOptionsToPackage{dvipsnames,svgnames,x11names}{xcolor}
%
\documentclass[
  letterpaper,
  DIV=11,
  numbers=noendperiod]{scrartcl}

\usepackage{amsmath,amssymb}
\usepackage{iftex}
\ifPDFTeX
  \usepackage[T1]{fontenc}
  \usepackage[utf8]{inputenc}
  \usepackage{textcomp} % provide euro and other symbols
\else % if luatex or xetex
  \usepackage{unicode-math}
  \defaultfontfeatures{Scale=MatchLowercase}
  \defaultfontfeatures[\rmfamily]{Ligatures=TeX,Scale=1}
\fi
\usepackage{lmodern}
\ifPDFTeX\else  
    % xetex/luatex font selection
\fi
% Use upquote if available, for straight quotes in verbatim environments
\IfFileExists{upquote.sty}{\usepackage{upquote}}{}
\IfFileExists{microtype.sty}{% use microtype if available
  \usepackage[]{microtype}
  \UseMicrotypeSet[protrusion]{basicmath} % disable protrusion for tt fonts
}{}
\makeatletter
\@ifundefined{KOMAClassName}{% if non-KOMA class
  \IfFileExists{parskip.sty}{%
    \usepackage{parskip}
  }{% else
    \setlength{\parindent}{0pt}
    \setlength{\parskip}{6pt plus 2pt minus 1pt}}
}{% if KOMA class
  \KOMAoptions{parskip=half}}
\makeatother
\usepackage{xcolor}
\setlength{\emergencystretch}{3em} % prevent overfull lines
\setcounter{secnumdepth}{-\maxdimen} % remove section numbering
% Make \paragraph and \subparagraph free-standing
\makeatletter
\ifx\paragraph\undefined\else
  \let\oldparagraph\paragraph
  \renewcommand{\paragraph}{
    \@ifstar
      \xxxParagraphStar
      \xxxParagraphNoStar
  }
  \newcommand{\xxxParagraphStar}[1]{\oldparagraph*{#1}\mbox{}}
  \newcommand{\xxxParagraphNoStar}[1]{\oldparagraph{#1}\mbox{}}
\fi
\ifx\subparagraph\undefined\else
  \let\oldsubparagraph\subparagraph
  \renewcommand{\subparagraph}{
    \@ifstar
      \xxxSubParagraphStar
      \xxxSubParagraphNoStar
  }
  \newcommand{\xxxSubParagraphStar}[1]{\oldsubparagraph*{#1}\mbox{}}
  \newcommand{\xxxSubParagraphNoStar}[1]{\oldsubparagraph{#1}\mbox{}}
\fi
\makeatother

\usepackage{color}
\usepackage{fancyvrb}
\newcommand{\VerbBar}{|}
\newcommand{\VERB}{\Verb[commandchars=\\\{\}]}
\DefineVerbatimEnvironment{Highlighting}{Verbatim}{commandchars=\\\{\}}
% Add ',fontsize=\small' for more characters per line
\usepackage{framed}
\definecolor{shadecolor}{RGB}{241,243,245}
\newenvironment{Shaded}{\begin{snugshade}}{\end{snugshade}}
\newcommand{\AlertTok}[1]{\textcolor[rgb]{0.68,0.00,0.00}{#1}}
\newcommand{\AnnotationTok}[1]{\textcolor[rgb]{0.37,0.37,0.37}{#1}}
\newcommand{\AttributeTok}[1]{\textcolor[rgb]{0.40,0.45,0.13}{#1}}
\newcommand{\BaseNTok}[1]{\textcolor[rgb]{0.68,0.00,0.00}{#1}}
\newcommand{\BuiltInTok}[1]{\textcolor[rgb]{0.00,0.23,0.31}{#1}}
\newcommand{\CharTok}[1]{\textcolor[rgb]{0.13,0.47,0.30}{#1}}
\newcommand{\CommentTok}[1]{\textcolor[rgb]{0.37,0.37,0.37}{#1}}
\newcommand{\CommentVarTok}[1]{\textcolor[rgb]{0.37,0.37,0.37}{\textit{#1}}}
\newcommand{\ConstantTok}[1]{\textcolor[rgb]{0.56,0.35,0.01}{#1}}
\newcommand{\ControlFlowTok}[1]{\textcolor[rgb]{0.00,0.23,0.31}{\textbf{#1}}}
\newcommand{\DataTypeTok}[1]{\textcolor[rgb]{0.68,0.00,0.00}{#1}}
\newcommand{\DecValTok}[1]{\textcolor[rgb]{0.68,0.00,0.00}{#1}}
\newcommand{\DocumentationTok}[1]{\textcolor[rgb]{0.37,0.37,0.37}{\textit{#1}}}
\newcommand{\ErrorTok}[1]{\textcolor[rgb]{0.68,0.00,0.00}{#1}}
\newcommand{\ExtensionTok}[1]{\textcolor[rgb]{0.00,0.23,0.31}{#1}}
\newcommand{\FloatTok}[1]{\textcolor[rgb]{0.68,0.00,0.00}{#1}}
\newcommand{\FunctionTok}[1]{\textcolor[rgb]{0.28,0.35,0.67}{#1}}
\newcommand{\ImportTok}[1]{\textcolor[rgb]{0.00,0.46,0.62}{#1}}
\newcommand{\InformationTok}[1]{\textcolor[rgb]{0.37,0.37,0.37}{#1}}
\newcommand{\KeywordTok}[1]{\textcolor[rgb]{0.00,0.23,0.31}{\textbf{#1}}}
\newcommand{\NormalTok}[1]{\textcolor[rgb]{0.00,0.23,0.31}{#1}}
\newcommand{\OperatorTok}[1]{\textcolor[rgb]{0.37,0.37,0.37}{#1}}
\newcommand{\OtherTok}[1]{\textcolor[rgb]{0.00,0.23,0.31}{#1}}
\newcommand{\PreprocessorTok}[1]{\textcolor[rgb]{0.68,0.00,0.00}{#1}}
\newcommand{\RegionMarkerTok}[1]{\textcolor[rgb]{0.00,0.23,0.31}{#1}}
\newcommand{\SpecialCharTok}[1]{\textcolor[rgb]{0.37,0.37,0.37}{#1}}
\newcommand{\SpecialStringTok}[1]{\textcolor[rgb]{0.13,0.47,0.30}{#1}}
\newcommand{\StringTok}[1]{\textcolor[rgb]{0.13,0.47,0.30}{#1}}
\newcommand{\VariableTok}[1]{\textcolor[rgb]{0.07,0.07,0.07}{#1}}
\newcommand{\VerbatimStringTok}[1]{\textcolor[rgb]{0.13,0.47,0.30}{#1}}
\newcommand{\WarningTok}[1]{\textcolor[rgb]{0.37,0.37,0.37}{\textit{#1}}}

\providecommand{\tightlist}{%
  \setlength{\itemsep}{0pt}\setlength{\parskip}{0pt}}\usepackage{longtable,booktabs,array}
\usepackage{calc} % for calculating minipage widths
% Correct order of tables after \paragraph or \subparagraph
\usepackage{etoolbox}
\makeatletter
\patchcmd\longtable{\par}{\if@noskipsec\mbox{}\fi\par}{}{}
\makeatother
% Allow footnotes in longtable head/foot
\IfFileExists{footnotehyper.sty}{\usepackage{footnotehyper}}{\usepackage{footnote}}
\makesavenoteenv{longtable}
\usepackage{graphicx}
\makeatletter
\newsavebox\pandoc@box
\newcommand*\pandocbounded[1]{% scales image to fit in text height/width
  \sbox\pandoc@box{#1}%
  \Gscale@div\@tempa{\textheight}{\dimexpr\ht\pandoc@box+\dp\pandoc@box\relax}%
  \Gscale@div\@tempb{\linewidth}{\wd\pandoc@box}%
  \ifdim\@tempb\p@<\@tempa\p@\let\@tempa\@tempb\fi% select the smaller of both
  \ifdim\@tempa\p@<\p@\scalebox{\@tempa}{\usebox\pandoc@box}%
  \else\usebox{\pandoc@box}%
  \fi%
}
% Set default figure placement to htbp
\def\fps@figure{htbp}
\makeatother

\KOMAoption{captions}{tableheading}
\makeatletter
\@ifpackageloaded{tcolorbox}{}{\usepackage[skins,breakable]{tcolorbox}}
\@ifpackageloaded{fontawesome5}{}{\usepackage{fontawesome5}}
\definecolor{quarto-callout-color}{HTML}{909090}
\definecolor{quarto-callout-note-color}{HTML}{0758E5}
\definecolor{quarto-callout-important-color}{HTML}{CC1914}
\definecolor{quarto-callout-warning-color}{HTML}{EB9113}
\definecolor{quarto-callout-tip-color}{HTML}{00A047}
\definecolor{quarto-callout-caution-color}{HTML}{FC5300}
\definecolor{quarto-callout-color-frame}{HTML}{acacac}
\definecolor{quarto-callout-note-color-frame}{HTML}{4582ec}
\definecolor{quarto-callout-important-color-frame}{HTML}{d9534f}
\definecolor{quarto-callout-warning-color-frame}{HTML}{f0ad4e}
\definecolor{quarto-callout-tip-color-frame}{HTML}{02b875}
\definecolor{quarto-callout-caution-color-frame}{HTML}{fd7e14}
\makeatother
\makeatletter
\@ifpackageloaded{caption}{}{\usepackage{caption}}
\AtBeginDocument{%
\ifdefined\contentsname
  \renewcommand*\contentsname{Table of contents}
\else
  \newcommand\contentsname{Table of contents}
\fi
\ifdefined\listfigurename
  \renewcommand*\listfigurename{List of Figures}
\else
  \newcommand\listfigurename{List of Figures}
\fi
\ifdefined\listtablename
  \renewcommand*\listtablename{List of Tables}
\else
  \newcommand\listtablename{List of Tables}
\fi
\ifdefined\figurename
  \renewcommand*\figurename{Figure}
\else
  \newcommand\figurename{Figure}
\fi
\ifdefined\tablename
  \renewcommand*\tablename{Table}
\else
  \newcommand\tablename{Table}
\fi
}
\@ifpackageloaded{float}{}{\usepackage{float}}
\floatstyle{ruled}
\@ifundefined{c@chapter}{\newfloat{codelisting}{h}{lop}}{\newfloat{codelisting}{h}{lop}[chapter]}
\floatname{codelisting}{Listing}
\newcommand*\listoflistings{\listof{codelisting}{List of Listings}}
\makeatother
\makeatletter
\makeatother
\makeatletter
\@ifpackageloaded{caption}{}{\usepackage{caption}}
\@ifpackageloaded{subcaption}{}{\usepackage{subcaption}}
\makeatother

\usepackage{bookmark}

\IfFileExists{xurl.sty}{\usepackage{xurl}}{} % add URL line breaks if available
\urlstyle{same} % disable monospaced font for URLs
\hypersetup{
  pdftitle={Hawaii SNAP Analysis: Interactive Homework},
  pdfauthor={DS 400: Bayesian Statistics \& Data Science},
  colorlinks=true,
  linkcolor={blue},
  filecolor={Maroon},
  citecolor={Blue},
  urlcolor={Blue},
  pdfcreator={LaTeX via pandoc}}


\title{Hawaii SNAP Analysis: Interactive Homework}
\usepackage{etoolbox}
\makeatletter
\providecommand{\subtitle}[1]{% add subtitle to \maketitle
  \apptocmd{\@title}{\par {\large #1 \par}}{}{}
}
\makeatother
\subtitle{Using Census Data \& Machine Learning to Understand Food
Security}
\author{DS 400: Bayesian Statistics \& Data Science}
\date{2025-10-15}

\begin{document}
\maketitle


\subsection{🌺 Introduction: Food Security in
Hawaiʻi}\label{introduction-food-security-in-hawaiux2bbi}

Food security is a critical health and equity issue. The Supplemental
Nutrition Assistance Program (SNAP) helps vulnerable households afford
nutritious food. But which communities have the highest rates of food
assistance participation? And what factors are most strongly associated
with SNAP participation?

In this assignment, you'll:

\begin{enumerate}
\def\labelenumi{\arabic{enumi}.}
\tightlist
\item
  \textbf{Access Census data} using \texttt{tidycensus}
\item
  \textbf{Explore patterns} of SNAP participation across Hawaiʻi census
  tracts
\item
  \textbf{Build a predictive model} using Random Forests
\item
  \textbf{Interpret results} using SHAP (Explainable AI) techniques
\end{enumerate}

\begin{tcolorbox}[enhanced jigsaw, title=\textcolor{quarto-callout-important-color}{\faExclamation}\hspace{0.5em}{Homework Requirements}, titlerule=0mm, coltitle=black, colbacktitle=quarto-callout-important-color!10!white, colback=white, left=2mm, toprule=.15mm, opacitybacktitle=0.6, toptitle=1mm, colframe=quarto-callout-important-color-frame, bottomtitle=1mm, arc=.35mm, bottomrule=.15mm, breakable, leftrule=.75mm, opacityback=0, rightrule=.15mm]

\begin{itemize}
\tightlist
\item
  ✅ Complete all \textbf{CHALLENGE} sections
\item
  ✅ Provide code AND written explanations
\item
  ✅ Show your thinking - comment your code!
\item
  ✅ Answer reflection questions at the end
\item
  ⭐ Bonus: Try the extension questions for extra credit
\end{itemize}

\textbf{Due Date}: {[}Your date here{]}

\end{tcolorbox}

\begin{center}\rule{0.5\linewidth}{0.5pt}\end{center}

\subsection{📦 Load Libraries}\label{load-libraries}

\begin{Shaded}
\begin{Highlighting}[]
\FunctionTok{library}\NormalTok{(tidycensus)}
\FunctionTok{library}\NormalTok{(dplyr)}
\FunctionTok{library}\NormalTok{(tidyr)}
\FunctionTok{library}\NormalTok{(stringr)}
\FunctionTok{library}\NormalTok{(mapgl)}
\FunctionTok{library}\NormalTok{(viridis)}
\FunctionTok{library}\NormalTok{(ggplot2)}
\FunctionTok{library}\NormalTok{(randomForest)}
\FunctionTok{library}\NormalTok{(treeshap)}
\FunctionTok{library}\NormalTok{(mapview)}
\FunctionTok{library}\NormalTok{(sf)}
\end{Highlighting}
\end{Shaded}

\begin{center}\rule{0.5\linewidth}{0.5pt}\end{center}

\subsection{📊 Part 1: Collect Census
Data}\label{part-1-collect-census-data}

\subsubsection{Set Up Your Census
Variables}\label{set-up-your-census-variables}

The code below defines the Census Bureau variables we'll use. These come
from the American Community Survey (ACS) 5-Year Estimates.

\begin{Shaded}
\begin{Highlighting}[]
\CommentTok{\# Define ACS variables for Hawaii}
\CommentTok{\# (Population, poverty, insurance, employment, housing, broadband, education, income, SNAP)}

\NormalTok{hi\_vars }\OtherTok{\textless{}{-}} \FunctionTok{c}\NormalTok{(}
  \CommentTok{\# Person{-}level}
  \AttributeTok{poverty\_n            =} \StringTok{"B17001\_002"}\NormalTok{,}
  \AttributeTok{uninsured\_u65\_n      =} \StringTok{"B27010\_017"}\NormalTok{,}
  \AttributeTok{unemployed\_n         =} \StringTok{"B23025\_005"}\NormalTok{,}
  \AttributeTok{pop\_total            =} \StringTok{"B01003\_001"}\NormalTok{,}
  
  \CommentTok{\# Households + housing}
  \AttributeTok{households\_total     =} \StringTok{"B11001\_001"}\NormalTok{,}
  \AttributeTok{broadband\_hh\_n       =} \StringTok{"B28002\_004"}\NormalTok{,}
  \AttributeTok{renter\_hh\_n          =} \StringTok{"B25003\_003"}\NormalTok{,}
  \AttributeTok{novehicle\_hh\_n       =} \StringTok{"B25044\_003"}\NormalTok{,}
  \AttributeTok{snap\_hh\_n            =} \StringTok{"B22001\_002"}\NormalTok{,     }\CommentTok{\# \textless{}{-}{-} SNAP households}
  
  \CommentTok{\# Education}
  \AttributeTok{edu\_25plus\_total     =} \StringTok{"B15003\_001"}\NormalTok{,}
  \AttributeTok{bachelors\_n          =} \StringTok{"B15003\_022"}\NormalTok{,}
  
  \CommentTok{\# Labor force + workers}
  \AttributeTok{labor\_force\_total    =} \StringTok{"B23025\_003"}\NormalTok{,}
  \AttributeTok{workers\_total        =} \StringTok{"B08301\_001"}\NormalTok{,}
  
  \CommentTok{\# Income level}
  \AttributeTok{median\_income        =} \StringTok{"B19013\_001"}
\NormalTok{)}
\end{Highlighting}
\end{Shaded}

\begin{tcolorbox}[enhanced jigsaw, title=\textcolor{quarto-callout-note-color}{\faInfo}\hspace{0.5em}{Understanding Census Variables}, titlerule=0mm, coltitle=black, colbacktitle=quarto-callout-note-color!10!white, colback=white, left=2mm, toprule=.15mm, opacitybacktitle=0.6, toptitle=1mm, colframe=quarto-callout-note-color-frame, bottomtitle=1mm, arc=.35mm, bottomrule=.15mm, breakable, leftrule=.75mm, opacityback=0, rightrule=.15mm]

Each variable code (like ``B17001\_002'') refers to a specific data
table and line number from the Census Bureau. The naming convention is:
- \textbf{B} = Detailed table - \textbf{17001} = Table ID (in this case,
poverty data) - \textbf{002} = Specific line number (usually the
numerator for a percentage calculation)

For more info, check the
\href{https://api.census.gov/data/2023/acs/acs5/variables.html}{Census
Bureau Data API documentation}

\end{tcolorbox}

\subsubsection{Pull Data from the Census
API}\label{pull-data-from-the-census-api}

\begin{Shaded}
\begin{Highlighting}[]
\CommentTok{\# Pull core variables (long format, then pivot to wide)}
\NormalTok{hi\_acs\_long }\OtherTok{\textless{}{-}} \FunctionTok{get\_acs}\NormalTok{(}
  \AttributeTok{geography =} \StringTok{"tract"}\NormalTok{,}
  \AttributeTok{state =} \StringTok{"HI"}\NormalTok{,}
  \AttributeTok{year =} \DecValTok{2023}\NormalTok{,}
  \AttributeTok{geometry =} \ConstantTok{TRUE}\NormalTok{,}
  \AttributeTok{variables =}\NormalTok{ hi\_vars}
\NormalTok{)}\SpecialCharTok{|\textgreater{}}
\NormalTok{  dplyr}\SpecialCharTok{::}\FunctionTok{filter}\NormalTok{(GEOID }\SpecialCharTok{!=} \StringTok{"15003981200"}\NormalTok{)}
\end{Highlighting}
\end{Shaded}

\begin{verbatim}

  |                                                                            
  |                                                                      |   0%
  |                                                                            
  |====                                                                  |   6%
  |                                                                            
  |=========                                                             |  13%
  |                                                                            
  |==========                                                            |  14%
  |                                                                            
  |============                                                          |  17%
  |                                                                            
  |=============                                                         |  19%
  |                                                                            
  |==============                                                        |  20%
  |                                                                            
  |===============                                                       |  22%
  |                                                                            
  |================                                                      |  23%
  |                                                                            
  |===================                                                   |  27%
  |                                                                            
  |====================                                                  |  29%
  |                                                                            
  |=====================                                                 |  30%
  |                                                                            
  |=======================                                               |  33%
  |                                                                            
  |=========================                                             |  35%
  |                                                                            
  |==========================                                            |  37%
  |                                                                            
  |===========================                                           |  38%
  |                                                                            
  |============================                                          |  40%
  |                                                                            
  |=============================                                         |  42%
  |                                                                            
  |===============================                                       |  45%
  |                                                                            
  |=================================                                     |  46%
  |                                                                            
  |==================================                                    |  48%
  |                                                                            
  |===================================                                   |  50%
  |                                                                            
  |====================================                                  |  51%
  |                                                                            
  |=====================================                                 |  53%
  |                                                                            
  |=======================================                               |  56%
  |                                                                            
  |=========================================                             |  58%
  |                                                                            
  |==========================================                            |  60%
  |                                                                            
  |===========================================                           |  61%
  |                                                                            
  |===============================================                       |  68%
  |                                                                            
  |====================================================                  |  74%
  |                                                                            
  |=======================================================               |  79%
  |                                                                            
  |=========================================================             |  81%
  |                                                                            
  |==========================================================            |  82%
  |                                                                            
  |===============================================================       |  90%
  |                                                                            
  |====================================================================  |  97%
  |                                                                            
  |======================================================================| 100%
\end{verbatim}

\begin{Shaded}
\begin{Highlighting}[]
\CommentTok{\# Convert to wide format (one row per census tract)}
\NormalTok{hi\_acs }\OtherTok{\textless{}{-}}\NormalTok{ hi\_acs\_long }\SpecialCharTok{\%\textgreater{}\%}
  \FunctionTok{select}\NormalTok{(GEOID, variable, estimate) }\SpecialCharTok{\%\textgreater{}\%}
  \FunctionTok{pivot\_wider}\NormalTok{(}\AttributeTok{names\_from =}\NormalTok{ variable, }\AttributeTok{values\_from =}\NormalTok{ estimate)}
\end{Highlighting}
\end{Shaded}

\subsubsection{Calculate Commute Times (≥30
minutes)}\label{calculate-commute-times-30-minutes}

\begin{Shaded}
\begin{Highlighting}[]
\CommentTok{\# Get commute time data using table lookup}
\NormalTok{commute\_tbl }\OtherTok{\textless{}{-}} \FunctionTok{get\_acs}\NormalTok{(}
  \AttributeTok{geography =} \StringTok{"tract"}\NormalTok{,}
  \AttributeTok{state =} \StringTok{"HI"}\NormalTok{,}
  \AttributeTok{year =} \DecValTok{2023}\NormalTok{,}
  \AttributeTok{table =} \StringTok{"B08303"}\NormalTok{,}
  \AttributeTok{geometry =} \ConstantTok{FALSE}
\NormalTok{)}\SpecialCharTok{|\textgreater{}}
\NormalTok{  dplyr}\SpecialCharTok{::}\FunctionTok{filter}\NormalTok{(GEOID }\SpecialCharTok{!=} \StringTok{"15003981200"}\NormalTok{)}

\CommentTok{\# Extract commute ≥30 min variable}
\NormalTok{commute\_30plus }\OtherTok{\textless{}{-}}\NormalTok{ commute\_tbl }\SpecialCharTok{\%\textgreater{}\%}
  \FunctionTok{filter}\NormalTok{(variable }\SpecialCharTok{==} \StringTok{"B08303\_004"}\NormalTok{) }\SpecialCharTok{\%\textgreater{}\%}
  \FunctionTok{group\_by}\NormalTok{(GEOID) }\SpecialCharTok{\%\textgreater{}\%}
  \FunctionTok{summarize}\NormalTok{(}\AttributeTok{commute\_30plus\_n =} \FunctionTok{sum}\NormalTok{(estimate, }\AttributeTok{na.rm =} \ConstantTok{TRUE}\NormalTok{), }\AttributeTok{.groups =} \StringTok{"drop"}\NormalTok{)}
\end{Highlighting}
\end{Shaded}

\begin{center}\rule{0.5\linewidth}{0.5pt}\end{center}

\subsection{🧮 Part 2: Calculate Percentages \& Clean
Data}\label{part-2-calculate-percentages-clean-data}

The code below joins all the data and calculates percentages. Study how
we normalize each variable!

\begin{Shaded}
\begin{Highlighting}[]
\CommentTok{\# Join commute data and calculate all percentages}
\NormalTok{hi\_acs }\OtherTok{\textless{}{-}}\NormalTok{ hi\_acs }\SpecialCharTok{\%\textgreater{}\%}
  \FunctionTok{left\_join}\NormalTok{(commute\_30plus, }\AttributeTok{by =} \StringTok{"GEOID"}\NormalTok{) }\SpecialCharTok{\%\textgreater{}\%}
  \FunctionTok{mutate}\NormalTok{(}
    \CommentTok{\# Person{-}based percentages}
    \AttributeTok{poverty\_rate\_pct        =} \DecValTok{100} \SpecialCharTok{*}\NormalTok{ poverty\_n       }\SpecialCharTok{/} \FunctionTok{pmax}\NormalTok{(pop\_total, }\DecValTok{1}\NormalTok{),}
    \AttributeTok{uninsured\_u65\_pct       =} \DecValTok{100} \SpecialCharTok{*}\NormalTok{ uninsured\_u65\_n }\SpecialCharTok{/} \FunctionTok{pmax}\NormalTok{(pop\_total, }\DecValTok{1}\NormalTok{),}
    
    \CommentTok{\# Labor{-}force based}
    \AttributeTok{unemployment\_rate\_pct   =} \DecValTok{100} \SpecialCharTok{*}\NormalTok{ unemployed\_n    }\SpecialCharTok{/} \FunctionTok{pmax}\NormalTok{(labor\_force\_total, }\DecValTok{1}\NormalTok{),}
    
    \CommentTok{\# Household{-}based percentages}
    \AttributeTok{broadband\_pct           =} \DecValTok{100} \SpecialCharTok{*}\NormalTok{ broadband\_hh\_n  }\SpecialCharTok{/} \FunctionTok{pmax}\NormalTok{(households\_total, }\DecValTok{1}\NormalTok{),}
    \AttributeTok{renter\_pct              =} \DecValTok{100} \SpecialCharTok{*}\NormalTok{ renter\_hh\_n     }\SpecialCharTok{/} \FunctionTok{pmax}\NormalTok{(households\_total, }\DecValTok{1}\NormalTok{),}
    \AttributeTok{no\_vehicle\_pct          =} \DecValTok{100} \SpecialCharTok{*}\NormalTok{ novehicle\_hh\_n  }\SpecialCharTok{/} \FunctionTok{pmax}\NormalTok{(households\_total, }\DecValTok{1}\NormalTok{),}
    \AttributeTok{snap\_pct                =} \DecValTok{100} \SpecialCharTok{*}\NormalTok{ snap\_hh\_n       }\SpecialCharTok{/} \FunctionTok{pmax}\NormalTok{(households\_total, }\DecValTok{1}\NormalTok{),}
    
    \CommentTok{\# Education (25+)}
    \AttributeTok{bachelors\_or\_higher\_pct =} \DecValTok{100} \SpecialCharTok{*}\NormalTok{ bachelors\_n     }\SpecialCharTok{/} \FunctionTok{pmax}\NormalTok{(edu\_25plus\_total, }\DecValTok{1}\NormalTok{),}
    
    \CommentTok{\# Commute{-}based}
    \AttributeTok{commute\_30plus\_pct      =} \DecValTok{100} \SpecialCharTok{*}\NormalTok{ commute\_30plus\_n }\SpecialCharTok{/} \FunctionTok{pmax}\NormalTok{(workers\_total, }\DecValTok{1}\NormalTok{)}
\NormalTok{  )}
\end{Highlighting}
\end{Shaded}

\begin{tcolorbox}[enhanced jigsaw, title=\textcolor{quarto-callout-tip-color}{\faLightbulb}\hspace{0.5em}{Why \texttt{pmax(denominator,\ 1)}?}, titlerule=0mm, coltitle=black, colbacktitle=quarto-callout-tip-color!10!white, colback=white, left=2mm, toprule=.15mm, opacitybacktitle=0.6, toptitle=1mm, colframe=quarto-callout-tip-color-frame, bottomtitle=1mm, arc=.35mm, bottomrule=.15mm, breakable, leftrule=.75mm, opacityback=0, rightrule=.15mm]

This prevents division by zero errors. If a census tract has 0 people,
we use 1 as the denominator instead. The \texttt{pmax()} function takes
the parallel maximum - comparing each value to 1 element-wise.

\end{tcolorbox}

\begin{center}\rule{0.5\linewidth}{0.5pt}\end{center}

\subsection{🗺️ CHALLENGE 1: Map SNAP
Participation}\label{challenge-1-map-snap-participation}

\begin{tcolorbox}[enhanced jigsaw, title=\textcolor{quarto-callout-warning-color}{\faExclamationTriangle}\hspace{0.5em}{Your Task}, titlerule=0mm, coltitle=black, colbacktitle=quarto-callout-warning-color!10!white, colback=white, left=2mm, toprule=.15mm, opacitybacktitle=0.6, toptitle=1mm, colframe=quarto-callout-warning-color-frame, bottomtitle=1mm, arc=.35mm, bottomrule=.15mm, breakable, leftrule=.75mm, opacityback=0, rightrule=.15mm]

Create a map showing SNAP participation rates by census tract across
Hawaiʻi. You should:

\begin{enumerate}
\def\labelenumi{\arabic{enumi}.}
\tightlist
\item
  \textbf{Use \texttt{ggplot2}} with \texttt{geom\_sf()} to map the data
\item
  \textbf{Color by \texttt{snap\_pct}} (use a viridis color scale)
\item
  \textbf{Add labels} including title, subtitle, and data source
\item
  \textbf{Make it look professional} with clean theme/formatting
\end{enumerate}

\textbf{Hints \& Resources:} - Use the \texttt{hi\_acs} dataset (it
already has geometry from \texttt{get\_acs()}) - Check out
\texttt{scale\_fill\_viridis()} for the color scale - Consider using
\texttt{color\ =\ NA} in \texttt{geom\_sf()} to remove tract borders -
Look at \texttt{theme\_minimal()} for a clean starting point - The code
structure might look like:
\texttt{r\ \ \ ggplot(hi\_acs)\ +\ \ \ \ \ geom\_sf(aes(fill\ =\ snap\_pct),\ color\ =\ NA)\ +\ \ \ \ \ scale\_fill\_viridis(option\ =\ "plasma",\ direction\ =\ -1,\ \ \ \ \ \ \ \ \ \ \ \ \ \ \ \ \ \ \ \ \ \ \ \ \ name\ =\ "\%\ with\ SNAP")\ +\ \ \ \ \ labs(title\ =\ "...",\ subtitle\ =\ "...",\ caption\ =\ "...")\ +\ \ \ \ \ theme\_minimal()\ +\ \ \ \ \ theme(...)}

\textbf{Your Code:}

\end{tcolorbox}

\begin{Shaded}
\begin{Highlighting}[]
\CommentTok{\# YOUR CODE HERE}
\CommentTok{\# Create a ggplot map of snap\_pct by census tract in Hawaii}
\CommentTok{\# Remember to include title, subtitle, and source citation!}
\end{Highlighting}
\end{Shaded}

\begin{tcolorbox}[enhanced jigsaw, title=\textcolor{quarto-callout-note-color}{\faInfo}\hspace{0.5em}{💭 Reflection on Your Map}, titlerule=0mm, coltitle=black, colbacktitle=quarto-callout-note-color!10!white, colback=white, left=2mm, toprule=.15mm, opacitybacktitle=0.6, toptitle=1mm, colframe=quarto-callout-note-color-frame, bottomtitle=1mm, arc=.35mm, bottomrule=.15mm, breakable, leftrule=.75mm, opacityback=0, rightrule=.15mm]

After creating your map, answer these questions in a comment:

\begin{enumerate}
\def\labelenumi{\arabic{enumi}.}
\tightlist
\item
  Where are the highest concentrations of SNAP participation?
\item
  Are there geographic patterns (islands, urban vs.~rural)?
\item
  What does this tell you about food security in Hawaiʻi?
\end{enumerate}

\end{tcolorbox}

\begin{center}\rule{0.5\linewidth}{0.5pt}\end{center}

\subsection{🔍 CHALLENGE 2: Interactive Map with
mapview}\label{challenge-2-interactive-map-with-mapview}

\begin{tcolorbox}[enhanced jigsaw, title=\textcolor{quarto-callout-warning-color}{\faExclamationTriangle}\hspace{0.5em}{Your Task}, titlerule=0mm, coltitle=black, colbacktitle=quarto-callout-warning-color!10!white, colback=white, left=2mm, toprule=.15mm, opacitybacktitle=0.6, toptitle=1mm, colframe=quarto-callout-warning-color-frame, bottomtitle=1mm, arc=.35mm, bottomrule=.15mm, breakable, leftrule=.75mm, opacityback=0, rightrule=.15mm]

Create an \textbf{interactive leaflet-style map} using the
\texttt{mapview} package.

\textbf{Hints \& Resources:} - The \texttt{mapview()} function creates
interactive maps - Use \texttt{zcol\ =\ "snap\_pct"} to color by SNAP
participation - Set \texttt{legend\ =\ TRUE} to show a legend - Use
\texttt{layer.name} to label your layer - The \texttt{at} parameter
controls color breaks (try: \texttt{seq(0,\ max(...),\ by\ =\ 5)}) - Use
\texttt{col.regions\ =\ viridis::viridis(10,\ direction\ =\ -1)} for
colors

\textbf{Your Code:}

\end{tcolorbox}

\begin{Shaded}
\begin{Highlighting}[]
\CommentTok{\# YOUR CODE HERE}
\CommentTok{\# Create an interactive mapview of snap\_pct}
\end{Highlighting}
\end{Shaded}

\begin{tcolorbox}[enhanced jigsaw, title=\textcolor{quarto-callout-tip-color}{\faLightbulb}\hspace{0.5em}{Interactive Map Features}, titlerule=0mm, coltitle=black, colbacktitle=quarto-callout-tip-color!10!white, colback=white, left=2mm, toprule=.15mm, opacitybacktitle=0.6, toptitle=1mm, colframe=quarto-callout-tip-color-frame, bottomtitle=1mm, arc=.35mm, bottomrule=.15mm, breakable, leftrule=.75mm, opacityback=0, rightrule=.15mm]

With \texttt{mapview}, you can: - Zoom and pan - Click on tracts to see
exact values - Toggle layers on/off - Switch between map types - Export
the map

Try zooming into a specific island and exploring the variation!

\end{tcolorbox}

\begin{center}\rule{0.5\linewidth}{0.5pt}\end{center}

\subsection{📊 CHALLENGE 3: Histogram of SNAP
Participation}\label{challenge-3-histogram-of-snap-participation}

\begin{tcolorbox}[enhanced jigsaw, title=\textcolor{quarto-callout-warning-color}{\faExclamationTriangle}\hspace{0.5em}{Your Task}, titlerule=0mm, coltitle=black, colbacktitle=quarto-callout-warning-color!10!white, colback=white, left=2mm, toprule=.15mm, opacitybacktitle=0.6, toptitle=1mm, colframe=quarto-callout-warning-color-frame, bottomtitle=1mm, arc=.35mm, bottomrule=.15mm, breakable, leftrule=.75mm, opacityback=0, rightrule=.15mm]

Create a histogram showing the distribution of \texttt{snap\_pct} across
all Hawaii census tracts.

\textbf{Hints \& Resources:} - Use \texttt{ggplot()} with
\texttt{geom\_histogram()} - Consider bin width - try
\texttt{bins\ =\ 20} or \texttt{binwidth\ =\ 2} - Add meaningful labels
- What does the shape tell you? Is it normally distributed?

\textbf{Your Code:}

\end{tcolorbox}

\begin{Shaded}
\begin{Highlighting}[]
\CommentTok{\# YOUR CODE HERE}
\CommentTok{\# Create a histogram of snap\_pct}
\end{Highlighting}
\end{Shaded}

\begin{tcolorbox}[enhanced jigsaw, title=\textcolor{quarto-callout-note-color}{\faInfo}\hspace{0.5em}{📈 Interpretation Questions}, titlerule=0mm, coltitle=black, colbacktitle=quarto-callout-note-color!10!white, colback=white, left=2mm, toprule=.15mm, opacitybacktitle=0.6, toptitle=1mm, colframe=quarto-callout-note-color-frame, bottomtitle=1mm, arc=.35mm, bottomrule=.15mm, breakable, leftrule=.75mm, opacityback=0, rightrule=.15mm]

\begin{enumerate}
\def\labelenumi{\arabic{enumi}.}
\tightlist
\item
  Is the distribution symmetric or skewed?
\item
  What's the typical SNAP participation rate?
\item
  Are there any outliers or surprising patterns?
\end{enumerate}

\end{tcolorbox}

\begin{center}\rule{0.5\linewidth}{0.5pt}\end{center}

\subsection{🌲 Part 3: Machine Learning
Setup}\label{part-3-machine-learning-setup}

Now we'll prepare for predictive modeling. First, create the ML-ready
dataset:

\begin{Shaded}
\begin{Highlighting}[]
\CommentTok{\# Create ML{-}ready dataframe with all our features}
\NormalTok{hi\_acs\_data\_ml }\OtherTok{\textless{}{-}}\NormalTok{ hi\_acs }\SpecialCharTok{\%\textgreater{}\%}
  \FunctionTok{select}\NormalTok{(GEOID, median\_income, poverty\_rate\_pct, uninsured\_u65\_pct, }
\NormalTok{         unemployment\_rate\_pct, broadband\_pct, renter\_pct, no\_vehicle\_pct, }
\NormalTok{         snap\_pct, bachelors\_or\_higher\_pct, commute\_30plus\_pct) }\SpecialCharTok{\%\textgreater{}\%}
  \FunctionTok{select}\NormalTok{(}\SpecialCharTok{{-}}\NormalTok{GEOID) }\SpecialCharTok{\%\textgreater{}\%}                              \CommentTok{\# Remove ID for modeling}
\NormalTok{  tidyr}\SpecialCharTok{::}\FunctionTok{drop\_na}\NormalTok{()                                }\CommentTok{\# Remove missing values}

\CommentTok{\# Check dimensions}
\FunctionTok{cat}\NormalTok{(}\StringTok{"Dataset dimensions:"}\NormalTok{, }\FunctionTok{dim}\NormalTok{(hi\_acs\_data\_ml), }\StringTok{"}\SpecialCharTok{\textbackslash{}n}\StringTok{"}\NormalTok{)}
\end{Highlighting}
\end{Shaded}

\begin{verbatim}
Dataset dimensions: 419 11 
\end{verbatim}

\begin{Shaded}
\begin{Highlighting}[]
\FunctionTok{cat}\NormalTok{(}\StringTok{"Features for modeling:"}\NormalTok{, }\FunctionTok{names}\NormalTok{(hi\_acs\_data\_ml))}
\end{Highlighting}
\end{Shaded}

\begin{verbatim}
Features for modeling: median_income poverty_rate_pct uninsured_u65_pct unemployment_rate_pct broadband_pct renter_pct no_vehicle_pct snap_pct bachelors_or_higher_pct commute_30plus_pct geometry
\end{verbatim}

\begin{center}\rule{0.5\linewidth}{0.5pt}\end{center}

\subsection{🎯 CHALLENGE 4: Build a Random Forest
Model}\label{challenge-4-build-a-random-forest-model}

\begin{tcolorbox}[enhanced jigsaw, title=\textcolor{quarto-callout-warning-color}{\faExclamationTriangle}\hspace{0.5em}{Your Task}, titlerule=0mm, coltitle=black, colbacktitle=quarto-callout-warning-color!10!white, colback=white, left=2mm, toprule=.15mm, opacitybacktitle=0.6, toptitle=1mm, colframe=quarto-callout-warning-color-frame, bottomtitle=1mm, arc=.35mm, bottomrule=.15mm, breakable, leftrule=.75mm, opacityback=0, rightrule=.15mm]

Train a Random Forest to predict SNAP participation rates using the
other socioeconomic indicators.

\textbf{Your Model Should:} - Use \texttt{snap\_pct} as the target
variable (left side of \texttt{\textasciitilde{}}) - Use all other
variables as predictors (\texttt{.} on right side) - Have
\texttt{ntree\ =\ 500} for stability - Be named \texttt{rf}

\textbf{Hints \& Resources:} - Use the \texttt{randomForest()} function
- Formula syntax: \texttt{target\ \textasciitilde{}\ .} predicts target
using all other columns - Structure:
\texttt{randomForest(formula,\ data\ =\ ...,\ ntree\ =\ 500)} - Print
the model to see performance metrics!

\textbf{Your Code:}

\end{tcolorbox}

\begin{Shaded}
\begin{Highlighting}[]
\CommentTok{\# YOUR CODE HERE}
\CommentTok{\# Build a random forest model to predict snap\_pct}
\CommentTok{\# Assign it to an object named \textquotesingle{}rf\textquotesingle{}}
\end{Highlighting}
\end{Shaded}

\begin{tcolorbox}[enhanced jigsaw, title=\textcolor{quarto-callout-note-color}{\faInfo}\hspace{0.5em}{🔍 Understanding Your Model Output}, titlerule=0mm, coltitle=black, colbacktitle=quarto-callout-note-color!10!white, colback=white, left=2mm, toprule=.15mm, opacitybacktitle=0.6, toptitle=1mm, colframe=quarto-callout-note-color-frame, bottomtitle=1mm, arc=.35mm, bottomrule=.15mm, breakable, leftrule=.75mm, opacityback=0, rightrule=.15mm]

Look at the printed output and answer:

\begin{enumerate}
\def\labelenumi{\arabic{enumi}.}
\tightlist
\item
  \textbf{Mean of squared residuals (MSR)}: How accurate are the
  predictions on average?
\item
  \textbf{\% Var explained}: What percentage of SNAP variation does the
  model explain?
\item
  \textbf{Interpretation}: Is this good? What might the model be
  missing?
\end{enumerate}

Remember: MSR tells you the average squared prediction error. \% Var
explained is like R² - higher is better!

\end{tcolorbox}

\begin{center}\rule{0.5\linewidth}{0.5pt}\end{center}

\subsection{🔬 Part 4: Explainable AI with
SHAP}\label{part-4-explainable-ai-with-shap}

SHAP (SHapley Additive exPlanations) values explain \emph{why} the model
makes specific predictions. They break down each prediction into
individual contributions from each feature.

\subsubsection{Understanding SHAP}\label{understanding-shap}

Before we calculate, let's understand what we're doing:

\begin{enumerate}
\def\labelenumi{\arabic{enumi}.}
\tightlist
\item
  \textbf{Unify the model}: Convert the Random Forest into a format that
  SHAP can work with
\item
  \textbf{Calculate SHAP values}: Compute explanations for a subset of
  observations (computing for all can be slow)
\item
  \textbf{Interpret}: Use the results to explain individual predictions
  and global patterns
\end{enumerate}

\begin{center}\rule{0.5\linewidth}{0.5pt}\end{center}

\subsection{🎯 CHALLENGE 8: Prepare SHAP
Calculations}\label{challenge-8-prepare-shap-calculations}

\begin{tcolorbox}[enhanced jigsaw, title=\textcolor{quarto-callout-warning-color}{\faExclamationTriangle}\hspace{0.5em}{Your Task}, titlerule=0mm, coltitle=black, colbacktitle=quarto-callout-warning-color!10!white, colback=white, left=2mm, toprule=.15mm, opacitybacktitle=0.6, toptitle=1mm, colframe=quarto-callout-warning-color-frame, bottomtitle=1mm, arc=.35mm, bottomrule=.15mm, breakable, leftrule=.75mm, opacityback=0, rightrule=.15mm]

Prepare the model for SHAP analysis by:

\begin{enumerate}
\def\labelenumi{\arabic{enumi}.}
\tightlist
\item
  \textbf{Unifying} your Random Forest model using \texttt{unify()}
\item
  \textbf{Calculating SHAP values} for the first 200 observations using
  \texttt{treeshap()}
\end{enumerate}

\textbf{Hints \& Resources:}

\begin{itemize}
\tightlist
\item
  Step 1: Use \texttt{unify(your\_model,\ your\_data)} to convert the RF

  \begin{itemize}
  \tightlist
  \item
    First argument: your Random Forest object (named \texttt{rf})
  \item
    Second argument: your ML dataset (named \texttt{hi\_acs\_data\_ml})
  \item
    Assign result to \texttt{unified}
  \end{itemize}
\item
  Step 2: Use
  \texttt{treeshap(unified\_object,\ data\_subset,\ verbose\ =\ 0)}

  \begin{itemize}
  \tightlist
  \item
    First argument: your unified model
  \item
    Second argument: subset of data - use
    \texttt{hi\_acs\_data\_ml{[}1:200,\ {]}} to get first 200 rows
  \item
    Third argument: \texttt{verbose\ =\ 0} suppresses progress messages
  \item
    Assign result to \texttt{treeshap1}
  \end{itemize}
\end{itemize}

\textbf{Your Code:}

\end{tcolorbox}

\begin{Shaded}
\begin{Highlighting}[]
\CommentTok{\# YOUR CODE HERE}
\CommentTok{\# Step 1: Unify your Random Forest model}
\CommentTok{\# unified \textless{}{-} unify(???, ???)}

\CommentTok{\# Step 2: Calculate SHAP values for first 200 observations}
\CommentTok{\# treeshap1 \textless{}{-} treeshap(???, ???, verbose = 0)}
\end{Highlighting}
\end{Shaded}

\begin{tcolorbox}[enhanced jigsaw, title=\textcolor{quarto-callout-note-color}{\faInfo}\hspace{0.5em}{💭 Why Only 200 Observations?}, titlerule=0mm, coltitle=black, colbacktitle=quarto-callout-note-color!10!white, colback=white, left=2mm, toprule=.15mm, opacitybacktitle=0.6, toptitle=1mm, colframe=quarto-callout-note-color-frame, bottomtitle=1mm, arc=.35mm, bottomrule=.15mm, breakable, leftrule=.75mm, opacityback=0, rightrule=.15mm]

Computing SHAP values is computationally expensive - it can take a long
time for large datasets. For this assignment, we're calculating for a
representative sample of 200 census tracts, which is enough to: -
Understand which features matter globally - Examine feature dependence
patterns - Look at individual case examples

In production, you might calculate for all observations or use
approximation methods like TreeSHAP (which is what we're using!).

\end{tcolorbox}

\begin{tcolorbox}[enhanced jigsaw, title=\textcolor{quarto-callout-tip-color}{\faLightbulb}\hspace{0.5em}{⚠️ Patience Required!}, titlerule=0mm, coltitle=black, colbacktitle=quarto-callout-tip-color!10!white, colback=white, left=2mm, toprule=.15mm, opacitybacktitle=0.6, toptitle=1mm, colframe=quarto-callout-tip-color-frame, bottomtitle=1mm, arc=.35mm, bottomrule=.15mm, breakable, leftrule=.75mm, opacityback=0, rightrule=.15mm]

The SHAP calculation might take 30 seconds to a few minutes depending on
your computer. The \texttt{verbose\ =\ 0} parameter keeps it quiet while
it works. You'll know it's done when the next code chunk starts
executing!

\end{tcolorbox}

\begin{center}\rule{0.5\linewidth}{0.5pt}\end{center}

\subsection{📊 CHALLENGE 5: Feature Importance
Plot}\label{challenge-5-feature-importance-plot}

\begin{tcolorbox}[enhanced jigsaw, title=\textcolor{quarto-callout-warning-color}{\faExclamationTriangle}\hspace{0.5em}{Your Task}, titlerule=0mm, coltitle=black, colbacktitle=quarto-callout-warning-color!10!white, colback=white, left=2mm, toprule=.15mm, opacitybacktitle=0.6, toptitle=1mm, colframe=quarto-callout-warning-color-frame, bottomtitle=1mm, arc=.35mm, bottomrule=.15mm, breakable, leftrule=.75mm, opacityback=0, rightrule=.15mm]

Create a feature importance plot showing the top 5 most important
variables for predicting SNAP participation.

\textbf{Hints \& Resources:} - Use
\texttt{plot\_feature\_importance(treeshap1,\ max\_vars\ =\ 5)} - This
uses the SHAP values we just calculated - The plot shows which features
have the biggest impact on predictions - Longer bars = more important

\textbf{Your Code:}

\end{tcolorbox}

\begin{Shaded}
\begin{Highlighting}[]
\CommentTok{\# YOUR CODE HERE}
\CommentTok{\# Plot feature importance (top 5 most important features)}
\end{Highlighting}
\end{Shaded}

\begin{tcolorbox}[enhanced jigsaw, title=\textcolor{quarto-callout-note-color}{\faInfo}\hspace{0.5em}{💭 Interpretation}, titlerule=0mm, coltitle=black, colbacktitle=quarto-callout-note-color!10!white, colback=white, left=2mm, toprule=.15mm, opacitybacktitle=0.6, toptitle=1mm, colframe=quarto-callout-note-color-frame, bottomtitle=1mm, arc=.35mm, bottomrule=.15mm, breakable, leftrule=.75mm, opacityback=0, rightrule=.15mm]

Based on your feature importance plot:

\begin{enumerate}
\def\labelenumi{\arabic{enumi}.}
\tightlist
\item
  What are the TOP 2 most important predictors of SNAP participation?
\item
  Why do you think these variables matter?
\item
  Are there any surprising results?
\item
  What does this tell you about food security in Hawaiʻi?
\end{enumerate}

\end{tcolorbox}

\begin{center}\rule{0.5\linewidth}{0.5pt}\end{center}

\subsection{📈 CHALLENGE 6: Feature Dependence
Plots}\label{challenge-6-feature-dependence-plots}

\begin{tcolorbox}[enhanced jigsaw, title=\textcolor{quarto-callout-warning-color}{\faExclamationTriangle}\hspace{0.5em}{Your Task}, titlerule=0mm, coltitle=black, colbacktitle=quarto-callout-warning-color!10!white, colback=white, left=2mm, toprule=.15mm, opacitybacktitle=0.6, toptitle=1mm, colframe=quarto-callout-warning-color-frame, bottomtitle=1mm, arc=.35mm, bottomrule=.15mm, breakable, leftrule=.75mm, opacityback=0, rightrule=.15mm]

Create feature dependence plots for the \textbf{top 5 most important
features} you identified in Challenge 5.

\textbf{What to Do:} 1. Look at your feature importance plot 2. Identify
the 5 most important variables 3. Create a dependence plot for each one
4. Use: \texttt{plot\_feature\_dependence(treeshap1,\ "variable\_name")}

\textbf{Example:}

\begin{Shaded}
\begin{Highlighting}[]
\FunctionTok{plot\_feature\_dependence}\NormalTok{(treeshap1, }\StringTok{"poverty\_rate\_pct"}\NormalTok{)}
\end{Highlighting}
\end{Shaded}

\textbf{Hints \& Resources:} - The x-axis shows the actual variable
values - The y-axis shows SHAP values (how much it pushes predictions
up/down) - Red dots = push toward higher SNAP \% - Blue dots = push
toward lower SNAP \% - Look for patterns: linear? non-linear? threshold
effects?

\textbf{Your Code (plot all 5):}

\end{tcolorbox}

\begin{Shaded}
\begin{Highlighting}[]
\CommentTok{\# YOUR CODE HERE {-} Feature Dependence Plot \#1}
\end{Highlighting}
\end{Shaded}

\begin{Shaded}
\begin{Highlighting}[]
\CommentTok{\# YOUR CODE HERE {-} Feature Dependence Plot \#2}
\end{Highlighting}
\end{Shaded}

\begin{Shaded}
\begin{Highlighting}[]
\CommentTok{\# YOUR CODE HERE {-} Feature Dependence Plot \#3}
\end{Highlighting}
\end{Shaded}

\begin{Shaded}
\begin{Highlighting}[]
\CommentTok{\# YOUR CODE HERE {-} Feature Dependence Plot \#4}
\end{Highlighting}
\end{Shaded}

\begin{Shaded}
\begin{Highlighting}[]
\CommentTok{\# YOUR CODE HERE {-} Feature Dependence Plot \#5}
\end{Highlighting}
\end{Shaded}

\begin{tcolorbox}[enhanced jigsaw, title=\textcolor{quarto-callout-note-color}{\faInfo}\hspace{0.5em}{📊 Dependence Analysis Questions}, titlerule=0mm, coltitle=black, colbacktitle=quarto-callout-note-color!10!white, colback=white, left=2mm, toprule=.15mm, opacitybacktitle=0.6, toptitle=1mm, colframe=quarto-callout-note-color-frame, bottomtitle=1mm, arc=.35mm, bottomrule=.15mm, breakable, leftrule=.75mm, opacityback=0, rightrule=.15mm]

For each plot, write 1-2 sentences about:

\begin{enumerate}
\def\labelenumi{\arabic{enumi}.}
\tightlist
\item
  \textbf{Direction}: Does higher variable value → higher or lower SNAP
  participation?
\item
  \textbf{Strength}: Is the relationship strong, moderate, or weak?
\item
  \textbf{Pattern}: Is it linear, or are there non-linear patterns?
\item
  \textbf{Outliers}: Are there interesting outlier tracts?
\end{enumerate}

Example answer: \textgreater{} ``Poverty rate shows a strong positive
relationship with SNAP participation. As poverty increases, SHAP values
increase substantially, pushing predictions toward higher SNAP
participation. The relationship appears roughly linear.''

\end{tcolorbox}

\begin{center}\rule{0.5\linewidth}{0.5pt}\end{center}

\subsection{👁️ CHALLENGE 7: Individual Tract
Contribution}\label{challenge-7-individual-tract-contribution}

\begin{tcolorbox}[enhanced jigsaw, title=\textcolor{quarto-callout-warning-color}{\faExclamationTriangle}\hspace{0.5em}{Your Task}, titlerule=0mm, coltitle=black, colbacktitle=quarto-callout-warning-color!10!white, colback=white, left=2mm, toprule=.15mm, opacitybacktitle=0.6, toptitle=1mm, colframe=quarto-callout-warning-color-frame, bottomtitle=1mm, arc=.35mm, bottomrule=.15mm, breakable, leftrule=.75mm, opacityback=0, rightrule=.15mm]

Examine what drives predictions for \textbf{Census Tract 304.04 on Maui,
Hawaiʻi}.

\textbf{Hints \& Resources:} - This tract is observation \#154 in our
data (hint: \texttt{obs\ =\ 154}) - Use:
\texttt{plot\_contribution(treeshap1,\ obs\ =\ 154)} - \textbf{Red bars}
= pushing toward higher SNAP \% - \textbf{Blue bars} = pushing toward
lower SNAP \% - \textbf{Bar length} = strength of effect

\textbf{Your Code:}

\end{tcolorbox}

\begin{Shaded}
\begin{Highlighting}[]
\CommentTok{\# YOUR CODE HERE}
\CommentTok{\# Plot contribution effects for Census Tract 304.04, Maui}
\end{Highlighting}
\end{Shaded}

\begin{tcolorbox}[enhanced jigsaw, title=\textcolor{quarto-callout-note-color}{\faInfo}\hspace{0.5em}{🔍 Case Study Analysis}, titlerule=0mm, coltitle=black, colbacktitle=quarto-callout-note-color!10!white, colback=white, left=2mm, toprule=.15mm, opacitybacktitle=0.6, toptitle=1mm, colframe=quarto-callout-note-color-frame, bottomtitle=1mm, arc=.35mm, bottomrule=.15mm, breakable, leftrule=.75mm, opacityback=0, rightrule=.15mm]

Answer these questions about Maui Census Tract 304.04:

\begin{enumerate}
\def\labelenumi{\arabic{enumi}.}
\tightlist
\item
  \textbf{Baseline}: What's the model's baseline prediction?
\item
  \textbf{Pushers}: Which 2-3 variables push predictions toward HIGHER
  SNAP\%?
\item
  \textbf{Pullers}: Which 2-3 variables push predictions toward LOWER
  SNAP\%?
\item
  \textbf{Net Effect}: Is the final prediction above or below the
  baseline?
\item
  \textbf{Story}: Write 2-3 sentences telling the ``story'' of this
  tract's SNAP participation
\end{enumerate}

\end{tcolorbox}

\begin{center}\rule{0.5\linewidth}{0.5pt}\end{center}

\subsection{🎓 Reflection \& Analysis}\label{reflection-analysis}

\begin{tcolorbox}[enhanced jigsaw, title=\textcolor{quarto-callout-important-color}{\faExclamation}\hspace{0.5em}{Final Reflection Questions}, titlerule=0mm, coltitle=black, colbacktitle=quarto-callout-important-color!10!white, colback=white, left=2mm, toprule=.15mm, opacitybacktitle=0.6, toptitle=1mm, colframe=quarto-callout-important-color-frame, bottomtitle=1mm, arc=.35mm, bottomrule=.15mm, breakable, leftrule=.75mm, opacityback=0, rightrule=.15mm]

Answer these questions in a separate document or as comments in your
code:

\subsubsection{Understanding the
Analysis}\label{understanding-the-analysis}

\begin{enumerate}
\def\labelenumi{\arabic{enumi}.}
\item
  \textbf{Model Performance}: How well does your Random Forest predict
  SNAP participation? What might explain the remaining unexplained
  variation?
\item
  \textbf{Key Drivers}: Based on feature importance and dependence
  plots, which 2-3 factors are most strongly associated with SNAP
  participation in Hawaiʻi?
\item
  \textbf{Equity Implications}:

  \begin{itemize}
  \tightlist
  \item
    What does the geographic distribution of SNAP participation tell you
    about food security disparities?
  \item
    Are certain communities more vulnerable?
  \end{itemize}
\item
  \textbf{Limitations}:

  \begin{itemize}
  \tightlist
  \item
    What important variables might be missing from this analysis?
  \item
    What are the limitations of predicting SNAP participation?
    (Remember: correlation ≠ causation)
  \end{itemize}
\item
  \textbf{Policy Recommendations}:

  \begin{itemize}
  \tightlist
  \item
    Based on your findings, what 2-3 policy interventions might help
    address food insecurity?
  \item
    Who should be targeted? Where should resources be focused?
  \end{itemize}
\end{enumerate}

\subsubsection{Technical Reflection}\label{technical-reflection}

\begin{enumerate}
\def\labelenumi{\arabic{enumi}.}
\setcounter{enumi}{5}
\item
  \textbf{Machine Learning Interpretation}: How did SHAP values help you
  understand the Random Forest model compared to just looking at
  variable importance?
\item
  \textbf{Next Steps}: If you were to continue this analysis, what would
  you do next?
\end{enumerate}

\end{tcolorbox}

\begin{center}\rule{0.5\linewidth}{0.5pt}\end{center}

\subsection{⭐ BONUS: Extension
Challenges}\label{bonus-extension-challenges}

\begin{tcolorbox}[enhanced jigsaw, title=\textcolor{quarto-callout-tip-color}{\faLightbulb}\hspace{0.5em}{For Extra Credit}, titlerule=0mm, coltitle=black, colbacktitle=quarto-callout-tip-color!10!white, colback=white, left=2mm, toprule=.15mm, opacitybacktitle=0.6, toptitle=1mm, colframe=quarto-callout-tip-color-frame, bottomtitle=1mm, arc=.35mm, bottomrule=.15mm, breakable, leftrule=.75mm, opacityback=0, rightrule=.15mm]

Try one or more of these extensions:

\subsubsection{Extension 1: Regional
Analysis}\label{extension-1-regional-analysis}

\begin{itemize}
\tightlist
\item
  Split Hawaiʻi by island (using \texttt{GEOID} or geographic location)
\item
  Build separate Random Forest models for each island
\item
  Compare which factors matter most on each island
\end{itemize}

\subsubsection{Extension 2: Relationship
Exploration}\label{extension-2-relationship-exploration}

\begin{itemize}
\tightlist
\item
  Create scatter plots showing relationships between top predictors and
  SNAP\%
\item
  Calculate correlation coefficients
\item
  Which relationships are strongest?
\end{itemize}

\subsubsection{Extension 3: Model
Comparison}\label{extension-3-model-comparison}

\begin{itemize}
\tightlist
\item
  Build a \textbf{second} Random Forest with different parameters (e.g.,
  \texttt{mtry}, \texttt{nodesize})
\item
  Compare performance to your original model
\item
  Which performs better?
\end{itemize}

\subsubsection{Extension 4: Prediction
Intervals}\label{extension-4-prediction-intervals}

\begin{itemize}
\tightlist
\item
  For a new hypothetical census tract, use your model to predict SNAP
  participation
\item
  What would be needed to create prediction intervals?
\item
  How confident would you be in that prediction?
\end{itemize}

\subsubsection{Extension 5: Data
Visualization}\label{extension-5-data-visualization}

\begin{itemize}
\tightlist
\item
  Create a multi-panel visualization combining:

  \begin{itemize}
  \tightlist
  \item
    The map of SNAP participation
  \item
    Distribution histogram
  \item
    Top feature importance plot
  \item
    Tell a complete visual story
  \end{itemize}
\end{itemize}

\end{tcolorbox}

\begin{center}\rule{0.5\linewidth}{0.5pt}\end{center}

\subsection{📝 Submission Checklist}\label{submission-checklist}

\begin{itemize}
\tightlist
\item[$\square$]
  All 7 challenges completed with code
\item[$\square$]
  Written explanations for each challenge
\item[$\square$]
  Reflection questions answered thoroughly
\item[$\square$]
  Code is commented and readable
\item[$\square$]
  Plots are properly labeled with titles/legends
\item[$\square$]
  At least one bonus challenge attempted (optional)
\end{itemize}

\begin{center}\rule{0.5\linewidth}{0.5pt}\end{center}

\subsection{📚 Resources \& References}\label{resources-references}

\begin{itemize}
\tightlist
\item
  \href{https://walker-data.com/tidycensus/}{tidycensus Documentation}
\item
  \href{https://data.census.gov/}{Census Bureau Data Catalog}
\item
  \href{https://shap.readthedocs.io/}{SHAP Documentation}
\item
  \href{https://cran.r-project.org/web/packages/randomForest/randomForest.pdf}{Random
  Forests in R}
\item
  \href{https://r-spatial.github.io/mapview/}{mapview Package}
\end{itemize}

\begin{center}\rule{0.5\linewidth}{0.5pt}\end{center}

\emph{Good luck! Remember: focus on understanding WHAT you're doing and
WHY, not just running code. Ask questions, explore the data, and enjoy
the detective work of data science! 🔍📊}




\end{document}
